\loai{Xác định các đại lượng theo công thức tính lực từ}
\begin{vidu} %[2P4-4.2B][Loại1]
Đặt một đoạn dây dẫn có chiều dài 2 m mang dòng điện 10 A vào một từ trường có cảm ứng từ là 0,02 T. Biết đường cảm ứng từ hợp với chiều dài của dây một góc là $60^\circ$. Lực từ tác dụng lên đoạn dây là
\choice
{0,3 N}
{0,519 N}
{\True 0,346 N}
{0,15 N}
\loigiai{Lực từ tác dụng lên dòng điện $F=IB\ell\sin \alpha =0,346\ N$.
}
\end{vidu}
\begin{vidu} %[2P4-4.2B][Loại1]
Một đoạn dây dẫn dài $\ell = 0,8\ m$ đặt trong từ trường đều sao cho dây dẫn hợp với vector cảm ứng từ một góc ${60^\circ}$. Biết dòng điện $I = 20\ A$ và dây dẫn chịu một lực là $F = 2.10^{-2}$ N. Độ lớn của cảm ứng từ là
\choice
{$0,8.10^{-3}$ T}
{$10^{-3}$ T}
{\True $1,4.10^{-3}$ T}
{$1,6.10^{-3}$ T}
\loigiai{Ta có:
$B=\dfrac{F}{I\ell \sin \alpha}=\dfrac{2.10^{-2}}{20.0,8.\sin 60^\circ}=1,44.10^{-3}\ T$}
\end{vidu}
\begin{vidu} %[2P4-4.2B][Loại1]
Một đoạn dây dẫn thẳng MN dài 6 cm có dòng điện 5 A đặt trong từ trường đều có cảm ứng từ $B = 0,5$ T. Lực từ tác dụng lên đoạn dây có độ lớn $7,5.10^{-2}$ N. Góc hợp bởi dây MN và đường cảm ứng từ là
\choice
{$90\degree$}
{$45\degree$}
{\True $30\degree$}
{$60\degree$}
\loigiai{Lực từ tác dụng lên đoạn dây $F = BI\ell\sin\alpha \Rightarrow \sin \alpha = \dfrac{F}{BIl} = \dfrac{7,5.10^{-2}}{0,5.5.0,06} = 0,5 \to \alpha = 30^\circ$.
}
\end{vidu}
\begin{vidu} %[2P4-4.2B][Loại1]
Một dây dẫn thẳng dài mang dòng điện 20 A, đặt trong từ trường đều có cảm ứng từ $B = 5.10^{-3}$ T. Dây dẫn đặt vuông góc với vector cảm ứng từ và chịu lực từ bằng $10^{-3}$ N. Chiều dài của đoạn dây dẫn là
\choice
{3 cm}
{\True 1 cm}
{2 cm}
{4 cm}
\loigiai{Lực từ tác dụng lên đoạn dây $F = BI\ell\sin\alpha \Rightarrow l = \dfrac{F}{BI\sin \alpha} = \dfrac{10^{-3}}{5.10^{-3}.20.\sin 90^\circ} = 0,01\ m = 1\ cm$}
\end{vidu}
\loai{Mối quan hệ giữa lực từ với cường độ dòng điện}
\begin{vidu} %[2P4-4.2B][Loại2]
Một đoạn dây dẫn mang dòng điện 4 A đặt trong một từ trường đều thì chịu một lực từ 8 N. Nếu dòng điện qua dây dẫn là 1 A thì nó chịu một lực từ có độ lớn bằng
\choice
{0,5 N}
{4 N}
{\True 2 N}
{32 N}
\loigiai{Lực từ tác dụng lên dây dẫn mang dòng điện đặt trong từ trường được xác định bằng biểu thức $.F=IB\ell\sin \alpha \Rightarrow $ Nếu dòng điện I giảm 4 lần thì F giảm 4 lần}
\end{vidu}
\loai{Lực từ tác dụng lên cạnh của khung dây}
\begin{vidu} %[2P4-4.2K][Loại3]
immini[thm]{Một dây dẫn được uốn thành một khung dây có dạng tam giác vuông tại A với $MN = 5$ cm, $AN = 3$ cm có dòng điện cường độ $I = 5$ A chạy qua. Đặt khung dây vào trong từ trường đều $B = 3. 10^{-3}$ T có vector cảm ứng từ song song với cạnh MN hướng như hình vẽ. Giữ khung dây cố định. Lực từ tác dụng lên cạnh AM có độ lớn
\choice
{$2,7. 10^{-4}$ N}
{$4,8. 10^{-3}$ N}
{\True $3,6. 10^{-4}$ N}
{$4,2. 10^{-3}$ N}}
{\begin{tikzpicture}[scale=0.7,thick,>=stealth']
\tkzDefPoints{0/0/a, 2/0/n, 0/3/m}
\draw[->]  (-3,3.5)  --++(-56:5);
\draw[->]  (-1,3.5)  --++(-56:5);
\draw[->]  (1,3.5)  --++(-56:5);
\tkzLabelPoint[left](a){$A$}
\tkzLabelPoint[right](n){$N$}
\tkzLabelPoint[left](m){$M$}
\tkzDrawSegments[arrow data={0.5}{stealth},very thick](a,m m,n n,a)
\node at (2.5,-1){$\vec{B}$};
\tkzLabelSegment[above](m,n){$I$}
\end{tikzpicture}}
\loigiai{Ta có: $\sin \alpha =\sin \widehat{AMN}\Rightarrow \sin\alpha =\dfrac{AN}{MN}=\dfrac{3}{5} \Rightarrow F=BI\ell \sin \alpha=3,6.10^{-4}\ N$}
\end{vidu}